\scanpage{27}
\uline{Formules d'\unsure{associativité|associativités}}

a) Vectoriels réels. \unsure{Soient|Soit} $I$ ens fini, $(E_i)_{i\in I}$ famille de vectoriels \inserted{réels} de dim finie, $E=\coprod_{i\in I}E_i$ la somme directe ($\simeq$ produit direct). \deleted{\unsure{Soit}} $I^{+}$ $(I^{-})$ l'ens des $i\in I$ tels que $d_i=\operatorname{rg}_{\mathbb R}E_i$ soit pair (impair). On \unsure{a un} iso\unsure{m} canonique
\[
 \boxed{\unsure{\omega|w}_E\simeq\omega_{I^{-}}\wedge\bigwedge_{i\in I}\omega_{E_i}}
\]
obtenu ainsi, en regardant le $2^{\text{ième}}$ mb comme \deleted{\unsure{un}} quotient de
{\reversemarginpar\margin{\figure{\begin{tikzpicture}
\node[inner sep=2pt] (nb) {\parbox{.86\linewidth}{\uline{NB} Un ordre total de $I$ en induit un de $I^{-}$, donc
\diagram{\resizebox{\linewidth}{!}{\begin{tikzcd}[ampersand replacement=\&,column sep=small]
\operatorname{numérot}(I)\arrow[r,"\text{épi}"] \& \operatorname{numérot}(I^{-})
\end{tikzcd}}}}};
\draw ([xshift=3pt]nb.north west) .. controls (nb.north) and ([xshift=-3pt]nb.north east) .. (nb.north east)
 .. controls ([xshift=-2pt]nb.east) and (nb.east) .. ([xshift=2pt]nb.south east)
 .. controls (nb.south east) and (nb.south) .. ([xshift=3pt]nb.south west);
\end{tikzpicture}}}\normalmarginpar}
\[
 \operatorname{numérot}(I^{-})\inserted{\times\operatorname{numérot}(I^{+})}\times\prod_{i\in I}\operatorname{bases}(E_i)
\]

Soit $\{i_1,\ldots,i_{n^{-}}\}$ \unsure{num.} de $I^{-}$, \inserted{$j_1\ \cdots\ j_{n^{+}}$ une autre de $I^{+}$,} et $\unsure{\forall}\ i\in I$, $B_{\unsure{i|\alpha}}=(\deleted{\unsure{B}}e_{i,1},\ldots,e_{i,d_i})$ une base ordonnée de $E_i$, alors
\[
 e_{i_1,1}\unsure{2|,},\ldots,e_{i_1,d_1},e_{i_2,\unsure{1|2}},\ldots e_{i_2,d_2},\ldots,e_{i_n,1},\ldots,e_{i_n,d_{i_n}}
\]
est une base de $E$. Si on change une $B_i$ par une matrice de $\mathrm{Gl}(d_i,\mathbb R)^{+}$, alors la base se trouve changée par une matrice de $\mathrm{Gl}(d,\mathbb R)^{+}$ $(d=\sum d_i\deleted{\unsure{)}}=\operatorname{rg}E)$, \inserted{donc définit \unsure{=|même} él. de $\omega_E$.} Si on change l'ordre de $I^{-}$ par une permutation $\sigma^{-}$ (une transposition entre $\alpha$ et $\alpha+1$, disons \ldots) On change la base par une matrice de dét $\deleted{\unsure{(-1)^{i_{\alpha}i_{\alpha+1}}|(-1)^{i_{\alpha}+i_{\alpha+1}}}}\operatorname{sgn}\sigma^{-}$ ; par contre, si on change \deleted{\unsure{par}} l'ordre de $I^{+}$ par une $\sigma^{+}$, la matrice de passage est de det $+1$ en tous cas \ldots
% archive sheet number: 26
